% Esempi di Tavolo TableTeX: gli stessi del pulsante «Esempi».
% Si compilano con pdflatex; aperto nell'editor, il file mostra le sue tabelle in una lista.
\documentclass{article}
\usepackage[T1]{fontenc}
\usepackage{amsmath,amssymb}
\usepackage[table]{xcolor}
\usepackage{booktabs,multirow,siunitx,tabularx,longtable,eurosym}

\begin{document}

\section*{Risultati con booktabs e siunitx}

\begin{table}[htbp]
  \centering
  \caption{Tempi di esecuzione (media di 10 prove)}
  \label{tab:tempi}
  \begin{tabular}{@{}l S[table-format=3.1] S[table-format=3.1] S[table-format=1.2]@{}}
    \toprule
    & \multicolumn{2}{c}{Tempo (\si{\milli\second})} & \\
    \cmidrule(lr){2-3}
    Algoritmo & {Medio} & {Peggiore} & {Rapporto} \\
    \midrule
    Quicksort  & 12.4  & 98.7  & 0.13 \\ % pivot casuale
    Mergesort  & 15.1  & 16.0  & 0.94 \\
    Heapsort   & 21.8  & 23.5  & 0.93 \\
    Timsort    & 9.6   & 15.2  & 0.63 \\ % implementazione di Python
    \bottomrule
  \end{tabular}
\end{table}

\section*{Griglia classica}

\begin{tabular}{|l|c|r|}
  \hline
  Prodotto & Quantità & Prezzo (\euro) \\
  \hline
  Quaderno & 3 & 4,50 \\
  \hline
  Penna    & 10 & 12,00 \\
  \hline
  Zaino    & 1 & 39,90 \\
  \hline
\end{tabular}

\section*{Celle unite}

\begin{tabular}{llcc}
  \toprule
  \multirow{2}{*}{Regione} & \multirow{2}{*}{Provincia} & \multicolumn{2}{c}{Popolazione} \\
  \cmidrule(lr){3-4}
  & & 2011 & 2021 \\
  \midrule
  \multirow{3}{*}{Campania} & Napoli  & 3054956 & 2988376 \\
                            & Salerno & 1092876 & 1070383 \\
                            & Caserta & 904921  & 916722  \\
  \midrule
  \multirow{2}{*}{Puglia}   & Bari    & 1247303 & 1226784 \\
                            & Lecce   & 802018  & 779090  \\
  \bottomrule
\end{tabular}

\section*{Testo lungo con tabularx}

\begin{tabularx}{\textwidth}{@{}lX@{}}
  \toprule
  Termine & Definizione \\
  \midrule
  Pivot & Elemento scelto per dividere il vettore in due parti: a sinistra i minori, a destra i maggiori. \\
  Stabilità & Un ordinamento è stabile se lascia nello stesso ordine gli elementi con chiavi uguali. \\
  In loco & L'algoritmo usa solo una quantità costante di memoria aggiuntiva. \\
  \bottomrule
\end{tabularx}

\section*{Elenco lungo con longtable}

\begin{longtable}{@{}rlr@{}}
  \caption{Elementi della tavola periodica}\label{tab:elementi} \\
  \toprule
  Z & Elemento & Massa (u) \\
  \midrule
  \endfirsthead
  \toprule
  Z & Elemento & Massa (u) \\
  \midrule
  \endhead
  \bottomrule
  \endfoot
  1 & Idrogeno & 1.008 \\
  2 & Elio & 4.003 \\
  3 & Litio & 6.94 \\
  4 & Berillio & 9.012 \\
  5 & Boro & 10.81 \\
  6 & Carbonio & 12.011 \\
  7 & Azoto & 14.007 \\
  8 & Ossigeno & 15.999 \\
\end{longtable}

\section*{Formule nelle celle}

\begin{tabular}{@{}lll@{}}
  \toprule
  Funzione & Derivata & Primitiva \\
  \midrule
  $x^n$ & $n x^{n-1}$ & $\frac{x^{n+1}}{n+1}$ \\
  $e^x$ & $e^x$ & $e^x$ \\
  $\ln x$ & $\frac{1}{x}$ & $x \ln x - x$ \\
  $\sin x$ & $\cos x$ & $-\cos x$ \\
  $\cos x$ & $-\sin x$ & $\sin x$ \\
  \bottomrule
\end{tabular}

\section*{Colori e intestazione in grassetto}

\begin{tabular}{lcc}
  \toprule
  \rowcolor{gray!15}
  \textbf{Squadra} & \textbf{Punti} & \textbf{Reti} \\
  \midrule
  Napoli & 82 & 77 \\
  \rowcolor{gray!8}
  Inter & 76 & 71 \\
  Milan & 70 & 64 \\
  \bottomrule
\end{tabular}

\end{document}
